
\subsubsection{\texorpdfstring{$\Delta$}{Delta}-Stepping 算法}

为了解决 Dijkstra 算法中严格的优先队列带来的难以并行化问题，Meyer 等人于 2003 年提出了 $\Delta$-Stepping 算法。该算法通过引入“桶（Bucket）”数据结构，在 Dijkstra 算法与 Bellman-Ford 算法之间取得了良好的平衡，从而实现了对图数据结构的有效并行处理。

算法思想：$\Delta$-Stepping 算法的核心在于将优先队列离散化为一系列连续的桶数组 $B$。设定一个步长参量 $\Delta > 0$，每个桶 $B_i$ 用于存放当前暂定距离 $d(v)$ 落在区间 $[i\Delta, (i+1)\Delta)$ 范围内的顶点。在每一轮迭代中，算法并发地从当前非空的编号最小的桶中取出所有顶点，并使用它们的临接“轻边”（权重 $\le \Delta$ 的边）去更新其后继顶点的距离。如果更新后的顶点距离落入了更小的桶区间，则将其移入对应的桶中。当当前桶内的顶点经过轻边完全松弛稳定后，再利用与之相连的“重边”（权重 $> \Delta$ 的边）进行一轮跨桶范围的并行扩展松弛。

时间复杂度：对于带最大度数限制的随机图拓扑结构，利用并行模型实现的 $\Delta$-Stepping 其平均时间复杂度可趋近于 $O(|V| + |E| + \frac{L}{\Delta})$，其中 $L$ 为最长最短路径距离。

在应对如线图高阶网格这种极度稠密网络且由于需要处理千万级节点状态的背景下，$\Delta$-Stepping 算法有效地避免了并行线程对于单一队列的锁争用（Lock Contention）。通过控制共享变量的并发原子访问（如使用 \texttt{atomicMin}），本算法可以在现代 GPU 计算管线上直接映射为无强制排序的桶内波前探索操作，这也是其在海量数据非加性路由重构系统内进行寻优加速的直接理论基础。

